\documentclass[../cours_tle.tex]{subfiles}
\begin{document}
\setcounter{chapter}{23}
\chapter{Algorithmes : calculer, prouver et contrôler l’erreur}
\label{ch:24}
\begin{objectif}
Spécialité. Exécuter les méthodes du cours et comprendre leurs limites. Prérequis : Python de première, chapitres 2, 11, 13 et 20.
\end{objectif}
\section{Variables, listes et fonctions}
Une affectation remplace la valeur d'une variable ; une égalité mathématique affirme une relation. Une fonction informatique possède des paramètres et renvoie un résultat ; préciser son contrat (entrées admises, sortie, erreur éventuelle).
Une liste Python est ordonnée, ses indices commencent à $0$ ; elle peut contenir des répétitions. Parcourir ses valeurs ou ses indices n'est pas la même opération.
\begin{lstlisting}[language=Python]
carres = [k*k for k in range(6)]
carres.append(36)
carres.pop()  # supprime le dernier element
pairs = [x for x in carres if x % 2 == 0]
\end{lstlisting}
\section{Dichotomie, Newton et sécante}
Pour $f$ continue avec $f(a)f(b)<0$, la dichotomie conserve l'invariant « une racine est dans $[a,b]$ et les signes aux bornes sont opposés ». La longueur est divisée par $2$ à chaque passage. Un nombre de passages fixé assure la terminaison et une borne d'erreur.

Newton remplace localement la courbe par sa tangente :
$x_{n+1}=x_n-f(x_n)/f'(x_n)$, si $f'(x_n)\ne0$.
La sécante remplace la dérivée par une pente entre deux points :
$x_{n+1}=x_n-f(x_n)(x_n-x_{n-1})/(f(x_n)-f(x_{n-1}))$.
Les dénominateurs doivent être contrôlés. Ces méthodes peuvent sortir du domaine, diverger ou cycler : pour $f(x)=x^3-2x+2$, Newton partant de $0$ alterne $0,1$.
Un petit résidu $|f(x)|$ ne garantit pas à lui seul une petite erreur sur $x$ ; si $|f'|\ge m>0$ entre $x$ et la racine, les accroissements finis donnent une erreur au plus $|f(x)|/m$.
\section{Euler et calcul approché d'intégrales}
Pour $y'=F(t,y)$, Euler explicite utilise $t_{k+1}=t_k+h$ et $y_{k+1}=y_k+hF(t_k,y_k)$.
Cela remplace l'évolution par sa tangente sur chaque petit intervalle. Pour $y'=ay+b$, c'est une suite affine ; l'erreur et la stabilité dépendent du pas.
Si $a<0$, l'erreur homogène est multipliée par $1+ah$ : un pas trop grand peut provoquer des oscillations divergentes alors que la solution exacte décroît.
Les rectangles, milieux et trapèzes du chapitre 20 calculent des intégrales par sommes ; leur preuve de convergence doit accompagner la formule.
\section{Approcher des constantes}
Héron approche $\sqrt2$ ; résoudre $x^2-x-1=0$ sur $[1,2]$ par dichotomie approche le nombre d'or. Les sommes $\sum_{k=0}^n1/k!$ approchent $e$ (première, chapitre 26).
L'aire d'un quart de disque unité donne $\pi=4\int_0^1\sqrt{1-x^2}\,dx$ ; des rectangles donnent un encadrement puisque la fonction décroît.
Briggs et Brouncker approchent $\ln2$ (chapitres 16 et 20).
Ces procédures ont des erreurs de troncature différentes et subissent aussi les arrondis informatiques.
\section{Monte-Carlo et échantillonnage}
Pour une fonction bornée sur $[a,b]$, tirer des $U_i$ uniformes indépendants puis calculer
$(b-a)n^{-1}\sum_i f(U_i)$ estime l'intégrale. L'identité $E(f(U))=(b-a)^{-1}\int_a^bf$ appartient au prolongement des probabilités aux densités ; elle est signalée ici comme fondement supplémentaire.
Pour l'aire d'un quart de disque, tirer $(x,y)$ uniformément dans le carré unité et compter $x^2+y^2\le1$. La proportion, multipliée par $4$, estime $\pi$ ; le chapitre 23 contrôle la fluctuation du comptage de Bernoulli.
Une graine rend une expérience reproductible, pas plus exacte.
\section{Les fichiers exécutables}
Le sous-dossier \texttt{python/} contient des fonctions pour Pascal, permutations, combinaisons, seuils, dichotomie, Newton, sécante, Euler, quadratures, Briggs, Brouncker, binomiale, Galton, Euclide et matrices.
Les tests comparent des cas exacts, des bornes d'erreur et des cas invalides. Ils ne constituent pas une preuve générale de la justesse mathématique.
En flottants, $0{,}1+0{,}2$ n'est pas représenté exactement par $0{,}3$ : pour une approximation, utiliser une tolérance argumentée ; pour un calcul rationnel exact, utiliser des fractions.
\begin{finchapitre}
\section{Utilité et suite du parcours}
La programmation permet d'expérimenter et d'évaluer le coût d'une méthode. En calcul scientifique, il faudra articuler preuve mathématique, erreur numérique et tests de logiciel.
\begin{retenir}
Un programme doit avoir une condition de terminaison et des entrées contrôlées. Une convergence observée n'est pas une garantie générale.
\end{retenir}
\end{finchapitre}
\end{document}
